%%%%%%%%%%%%%%%%%%%%%%%%%%%%%%%%%%%%%%%%%%%%%%%%%%%%%%%%%%%%%%%%%%%%%%%%%%%%%%%%
\section{Baseline Predictive Performance on Full AIGs (RQ1)}
\label{sec:results:rq1}
% Answers \ref{sec:intro:rq1}. Reference point for every reduction result downstream.

\subsection{Comparison Against Naive Baselines}
\label{sec:results:rq1:baselines}
``Accurately'' needs a referent. Two groups supply one
(\ref{sec:method:architecture:baselines}): trivial predictors, which fix what a given $R^2$
is worth on this label, and published circuit models ported to this task.

\begin{figure*}[htbp]
    \centering
    \includegraphics[width=\textwidth]{\resultsfigs/rq1_baseline_tiers}
    \caption[Baseline group comparison]{\textbf{PARTLY FABRICATED.} Every RQ1 baseline on
    the design-disjoint test split, by RMSE, $R^2$ and Spearman correlation. The trivial
    predictors, fitted on validation and scored on test, and the edge-aware GCN+ are
    measured; the published circuit models are invented placeholders. The size-only
    regressor outranks the encoder on Spearman correlation.}
    \label{fig:rq1_baseline_tiers}
\end{figure*}

The grouped comparison in~\ref{fig:rq1_baseline_tiers} is the reference point for the whole
chapter: it fixes what the encoder is worth against predictors that read nothing but the
label distribution or the node and edge counts. The size-only ordinary least squares fit
already ranks circuits better than the encoder does, so a claim that the encoder reads
circuit \emph{structure} rather than circuit \emph{scale} has to clear that row before any
reduction result means anything.
\clearpage

\subsection{Predictive Accuracy Metrics}
\label{sec:results:rq1:accuracy}
Smooth L1 loss, RMSE and $R^2$ for the edge-aware GCN+ predicting Orchestrate
optimizability on a test set of full AIGs, reported together
(\ref{sec:method:experiment:metrics}).

\begin{figure}[htbp]
    \centering
    \includegraphics[width=\columnwidth]{\resultsfigs/rq1_parity}
    \caption[Full-graph parity plot]{Predicted against true optimizability for the
    full-graph baseline on the design-disjoint test split, one point per circuit. The
    diagonal is perfect prediction: points below it are under-predictions, points above it
    over-predictions.}
    \label{fig:rq1_parity}
\end{figure}

A parity plot is the standard first look at a regression fit, and it is here because two
failures that no single error number exposes are visible in it directly. A model that has
compressed toward the label mean shows as a cloud flattened onto a horizontal band instead
of following the diagonal, however healthy its RMSE looks. A model with heteroscedastic
error shows as a spread about the diagonal that widens with the target, which says the
reported error is an average over regimes that behave differently rather than a property of
the model as a whole. Both matter for a label that is 49\% exactly zero.

\begin{figure}[htbp]
    \centering
    \includegraphics[width=\columnwidth]{\resultsfigs/rq1_calibration}
    \caption[Baseline calibration]{Binned reliability of the full-graph baseline: mean
    prediction against mean truth within bins of predicted optimizability. A model that has
    collapsed toward the label mean shows here as a flat curve regardless of how healthy its
    RMSE looks.}
    \label{fig:rq1_calibration}
\end{figure}

Calibration takes the same question as~\ref{fig:rq1_parity} and answers it by bin rather
than by point, which is the readable form when most of the mass sits at one end of the
range.

\begin{figure*}[htbp]
    \centering
    \includegraphics[width=\textwidth]{\resultsfigs/rq1_error_by_target_bin}
    \caption[Baseline performance by optimizability band]{The full-graph baseline as a
    function of how much the circuit actually needs optimizing. The model tracks the truth
    on the two lowest bands, which hold 80\% of the corpus, and then flattens: at a true
    optimizability above 0.20 it predicts 0.046 against a true mean of 0.226, an absolute
    error of 0.180. It has largely learned the label's marginal distribution rather than the
    task.}
    \label{fig:rq1_error_by_target_bin}
\end{figure*}

The banded view in~\ref{fig:rq1_error_by_target_bin} is where the flattening
of~\ref{fig:rq1_calibration} becomes a number. The circuits the motivating use case cares
about most, the ones with the most to remove, are exactly the ones the baseline predicts
worst.

\begin{figure*}[htbp]
    \centering
    \includegraphics[width=\textwidth]{\resultsfigs/rq1_variance_decomposition}
    \caption[Variance decomposition of the baseline score]{How much of the full-graph
    baseline's $R^2 = 0.343$ is reachable without learning anything about an individual
    circuit. Grouping by dataset tier carries only 1.1\% of the label variance, so the tier
    is not a confound. Grouping by \emph{source script} carries 25.0\%: a predictor that
    outputs nothing but each group's mean scores $R^2 = 0.250$, roughly three-quarters of
    the model's reported figure. Held against within-group variance the model retains
    $R^2 = 0.124$.}
    \label{fig:rq1_variance_decomposition}
\end{figure*}

The decomposition in~\ref{fig:rq1_variance_decomposition} separates the headline score into
the part any group-mean predictor would get for free and the part attributable to reading
circuit structure. The residual $R^2 = 0.124$ is the quantity this section defends, and it
is the figure every reduction in RQ3 should be compared against rather than the pooled
0.343.
\clearpage

\subsection{Ranking Efficacy}
\label{sec:results:rq1:ranking}
Spearman correlation between predicted and true optimizability, which is the metric the
script-selection use case in~\ref{sec:intro:motivation} actually depends on.

\begin{figure}[htbp]
    \centering
    \fbox{\parbox[c][0.24\textheight][c]{0.86\columnwidth}{\centering\footnotesize
    Reserved slot. No plot generated.}}
    \caption[Ranking efficacy of the full-graph baseline]{\textbf{FABRICATED.} Reserved slot
    for the ranking view of the RQ1 comparison: Spearman correlation between predicted and
    true optimizability for each baseline on the design-disjoint test split, with a
    bootstrap interval per bar. Nothing is plotted and no data stands behind this float. The
    numbers it would draw are the Spearman column of~\ref{tab:rq1_baselines}.}
    \label{fig:rq1_ranking}
\end{figure}

Ranking is reported separately from error throughout because the two can move in opposite
directions, and~\ref{fig:rq1_baseline_tiers} already shows them doing so: the encoder wins
on RMSE and $R^2$ and loses on Spearman correlation to a two-parameter fit on graph size.
\clearpage

\subsection{Baseline Model Outcomes}
\label{sec:results:rq1:baselinemodels}
The published circuit models were adapted rather than reused, and the adaptations are stated
in~\ref{sec:method:architecture:baselines}. What each comparison licenses depends on them.

\begin{table*}[htbp]
\centering
\caption{RQ1 three-tier baseline comparison on the design-disjoint test split. Trivial predictors are fitted on validation and scored on test. Upstream $R^2$ figures are NOT comparable: OpenABC-D z-scores its targets per design, which removes between-design variance from the denominator. Only the sign transfers.}
\label{tab:rq1_baselines}
\small
\begin{tabular}{llrrrr}
\toprule
\textbf{Model} & \textbf{Tier} & \textbf{RMSE} & \textbf{MAE} & \textbf{$R^2$} & \textbf{Spearman} \\
\midrule
Predict validation mean & 1: trivial & 0.054 & 0.023 & -0.052 & -- \\
Predict validation median & 1: trivial & 0.056 & 0.020 & -0.131 & -- \\
Graph size only (OLS) & 1: trivial & 0.050 & 0.024 & 0.101 & 0.348 \\
Edge-aware GCN+ (ours) & This work & 0.043 & 0.017 & 0.343 & 0.185 \\
{[TODO/FAKE] GCN (mean pool)} & 2: standard encoders & 0.050 & 0.030 & 0.100 & 0.150 \\
{[TODO/FAKE] GraphSAGE} & 2: standard encoders & 0.048 & 0.029 & 0.150 & 0.200 \\
{[TODO/FAKE] GIN} & 2: standard encoders & 0.046 & 0.028 & 0.200 & 0.250 \\
{[TODO/FAKE] SynthNet} & 3: published models & 0.021 & 0.012 & -0.168 & 0.000 \\
{[TODO/FAKE] HOGA} & 3: published models & 0.050 & 0.030 & 0.100 & 0.100 \\
{[TODO/FAKE] DeepGate4} & 3: published models & 0.045 & 0.027 & 0.250 & 0.300 \\
\bottomrule
\end{tabular}
\\[4pt]\begin{minipage}{\linewidth}\footnotesize\emph{Rows marked [TODO/FAKE] are invented. Train GCN / GraphSAGE / GIN under the headline pooling, head and budget (src/shell/train.sh with --encoder\_name), and finish the DeepGate4 and HOGA ports on the baselines/openabc-synthnet-hoga branch. SynthNet is the only published baseline that has produced a number (train\_baseline\_synthnet\_Orchestrate, val R2 = -0.168); the HOGA run crashed at epoch 0.}\end{minipage}
\end{table*}


The table collects all of RQ1 in one place, and its footnote records which rows are invented
and which measured. Two cautions attach to it. SynthNet is the only published baseline that
has produced a number at all, and it collapsed to a constant prediction on this split, so
its RMSE looks strong while its $R^2$ is negative. Its published $R^2$ is not the same
quantity as the one reported here, because the upstream work z-scores its targets per
design and so removes between-design variance from the denominator; only the sign transfers.
\clearpage

\subsection{Error Analysis}
\label{sec:results:rq1:erroranalysis}
Residuals broken down by graph size, by tier and by source design, which sets the
expectation for which reduction methods should be expected to hurt most.

\begin{figure*}[htbp]
    \centering
    \includegraphics[width=\textwidth]{\resultsfigs/rq1_residuals_by_size}
    \caption[Residuals by circuit scale]{Absolute error of the full-graph baseline against
    graph size (median and interquartile band), with the number of graphs per size bin
    alongside so that a wide band on a thin bin is not read as a trend.}
    \label{fig:rq1_residuals_by_size}
\end{figure*}

Whether error grows with scale is the precondition for the reduction chapters: a method that
acts hardest on large graphs starts from whatever position~\ref{fig:rq1_residuals_by_size}
shows, not from the pooled error.

\begin{figure}[htbp]
    \centering
    \includegraphics[width=\columnwidth]{\resultsfigs/rq1_error_by_tier}
    \caption[Error by optimization tier]{Absolute error of the full-graph baseline by
    dataset tier. Tiers differ by construction, since a graph already optimized once has
    less left to remove.}
    \label{fig:rq1_error_by_tier}
\end{figure}

A difference across tiers is expected, and its absence would itself be a finding about the
generation pipeline rather than about the model. Against
\ref{fig:rq1_variance_decomposition}, which puts only 1.1\% of the label variance on the
tier, the tier is not the axis that structures this corpus.

\begin{figure*}[htbp]
    \centering
    \includegraphics[width=\textwidth]{\resultsfigs/rq1_error_by_design}
    \caption[Per-design error on unseen designs]{Per-design RMSE and $R^2$ for the
    full-graph baseline on the design-disjoint test set. The spread across designs is very
    wide and $R^2$ is strongly negative on some of them.}
    \label{fig:rq1_error_by_design}
\end{figure*}

\begin{table*}[htbp]
\centering
\caption{Per-design metrics on the design-disjoint test set. A pooled score over a design-level split is an average over this handful of whole circuits, not over a random sample.}
\label{tab:rq1_per_design}
\small
\begin{tabular}{lrrrrrr}
\toprule
\textbf{Design} & \textbf{Graphs} & \textbf{Mean $y$} & \textbf{RMSE} & \textbf{MAE} & \textbf{$R^2$} & \textbf{Spearman} \\
\midrule
256 & 16040 & 0.001 & 0.001 & 0.001 & -0.548 & 0.493 \\
16384 & 16080 & 0.000 & 0.001 & 0.001 & -747.953 & 0.043 \\
picosoc & 16045 & 0.007 & 0.010 & 0.006 & -0.657 & 0.265 \\
c1355 & 16102 & 0.006 & 0.026 & 0.006 & -0.014 & 0.331 \\
c6288 & 16052 & 0.044 & 0.055 & 0.038 & 0.347 & 0.539 \\
square & 16111 & 0.062 & 0.085 & 0.049 & 0.091 & 0.884 \\
\bottomrule
\end{tabular}
\end{table*}


The figure and the table give the same result in two forms, and it is the most consequential
result in RQ1: a pooled score over a design-level split is an average over a handful of
whole circuits, not over a random sample, so the pooled number describes the mix of designs
at least as much as it describes the model. That is the observation RQ1a
(\ref{sec:results:rq1:protocol}) turns into a measurement.
\clearpage

\subsection{Protocol Sensitivity: How Much of the Score Is Design Recognition? (RQ1a)}
\label{sec:results:rq1:protocol}
Three train and test protocols under otherwise identical conditions: a plain random split, a
recipe-disjoint split, and the design-disjoint split used for every reported result
(\ref{sec:method:experiment:splitting}).

\begin{figure*}[htbp]
    \centering
    \includegraphics[width=\textwidth]{\resultsfigs/rq1a_split_protocol}
    \caption[Split protocol sensitivity]{\textbf{MOSTLY FABRICATED.} How much of the
    measured score is design recognition rather than structural inference. Only the
    design-disjoint row is measured; the random and recipe-disjoint rows are placeholders
    for two training runs that have not completed.}
    \label{fig:rq1a_split_protocol}
\end{figure*}

\begin{table}[htbp]
\centering
\caption{RQ1a protocol sensitivity. Identical encoder, budget and seed; only the split changes. Random is the common default in circuit representation-learning evaluations, recipe-disjoint is OpenABC-D Variant~1, design-disjoint is Variant~2 and is what this thesis reports.}
\label{tab:rq1a_protocol}
\small
\begin{tabular}{lrrrr}
\toprule
\textbf{Protocol} & \textbf{RMSE} & \textbf{$R^2$} & \textbf{Spearman} & \textbf{Inflation vs. design-disjoint} \\
\midrule
\textcolor{red}{{[TODO/FAKE] Random (per-row)}} & \textcolor{red}{999999} & \textcolor{red}{-999999} & \textcolor{red}{-999999} & \textcolor{red}{999999} \\
\textcolor{red}{{[TODO/FAKE] Recipe-disjoint}} & \textcolor{red}{999999} & \textcolor{red}{-999999} & \textcolor{red}{-999999} & \textcolor{red}{999999} \\
Design-disjoint (ours) & 0.043 & 0.343 & 0.185 & 0.000 \\
\bottomrule
\end{tabular}
\\[4pt]\begin{minipage}{\linewidth}\footnotesize\emph{\textcolor{red}{Rows in red are FABRICATED PLACEHOLDERS, not measurements: every number in them has been replaced by the sentinel $\pm$999999.} Rows marked [TODO/FAKE] are invented. IN FLIGHT: the --split\_by random and --split\_by recipe training runs are on the cluster now (the design-disjoint run already exists). Identical encoder, budget and seed; only the split changes. When they land, run src/test.py for each, delete this block and rerun analysis.make\_all.}\end{minipage}
\end{table}


The gap between the three rows is a direct measurement of how much apparent accuracy a
permissive protocol buys, and until the two outstanding runs land, the comparison against
published numbers obtained under leakier protocols cannot be made quantitative. The
placeholder rows are set to plausible values only so that the figure has a shape; no
conclusion may be drawn from their size.
\clearpage

\subsection{Cost of the Full-Graph Baseline}
\label{sec:results:rq1:hardware}
Three quantities characterise the cost of training on unreduced graphs, and together they
form the reference against which every reduction in RQ2 is read: time per optimisation step
and per epoch, peak allocated device memory, and inference throughput in graphs per second.

\begin{figure*}[htbp]
    \centering
    \includegraphics[width=\textwidth]{\resultsfigs/rq1_training_curves}
    \caption[Training trajectories]{Validation loss and $R^2$ per epoch for every
    Orchestrate training run, coloured by reduction family. Validation $R^2$ swings by
    several units between consecutive epochs on some configurations.}
    \label{fig:rq1_training_curves}
\end{figure*}

The swing bounds how much weight any single-seed comparison can carry
(\ref{sec:method:experiment:metrics:stats}). Every configuration in this thesis is trained
exactly once, so a gap of a few percent between two configurations cannot be separated from
the epoch-to-epoch movement visible here.

\begin{figure*}[htbp]
    \centering
    \includegraphics[width=\textwidth]{\resultsfigs/rq1_training_cost}
    \caption[Training budget consumed]{Wall-clock time and epochs actually consumed per
    configuration, against epochs to best validation loss. Two runs of the same nominal
    budget are not the same amount of compute if one stopped at epoch 4 and the other at
    10.}
    \label{fig:rq1_training_cost}
\end{figure*}

Consumed budget rather than allotted budget is what the RQ2 speedups are normalised
against, since early stopping makes the two diverge.

\begin{figure*}[htbp]
    \centering
    \includegraphics[width=\textwidth]{\resultsfigs/rq1_hp_sweep}
    \caption[Hyperparameter sweep]{\textbf{APPENDIX.} The hyperparameter sweep that fixed
    the architecture: progress per worker, objective by encoder, and trial outcomes. A large
    share of trials were pruned on memory exhaustion rather than on their objective.
    Reconstructed from the worker logs, because the study databases were lost with the
    scratch workspaces.}
    \label{fig:rq1_hp_sweep}
\end{figure*}

Memory pressure being the dominant failure mode of the search is the hyperparameter-search
face of the bottleneck this thesis is about, which is the one reason the sweep is reported
in the body at all.

\begin{figure*}[htbp]
    \centering
    \includegraphics[width=\textwidth]{\resultsfigs/rq1_hp_parameters}
    \caption[Hyperparameter sensitivity]{\textbf{APPENDIX.} Objective against each searched
    hyperparameter over the 52 completed trials. A coarse sensitivity read rather than an
    importance analysis: sampling concentrates on promising regions and the sampler state
    was lost with the scratch workspaces.}
    \label{fig:rq1_hp_parameters}
\end{figure*}

\begin{table}[htbp]
\centering
\caption{Final configuration, as recorded in the run log of the reported run. Fields the run did not log are marked rather than omitted.}
\label{tab:rq1_hyperparameters}
\small
\begin{tabular}{ll}
\toprule
\textbf{Hyperparameter} & \textbf{Value} \\
\midrule
Encoder & gcn \\
Hidden dim & 128.0 \\
Layers (message passing) & 4 (not logged) \\
Positional encoding & level \\
PE dim & 32.0 \\
Pooling & mean \\
Head dropout & 0.3 \\
Learning rate & 0.0003 \\
Min learning rate & 1e-06 \\
Weight decay & 0.00396 \\
Warmup steps & 10000.0 \\
Scheduler patience & 2.0 \\
Node input dim & 4.0 \\
Edge attr dim & 2.0 \\
Seed & not recorded \\
\bottomrule
\end{tabular}
\end{table}


The sensitivity view supports ``which settings never produced a good trial'' and not ``which
parameter matters most''. The table records the final configuration as logged, marking the
fields the run did not log rather than omitting them; it moves to the appendix with the two
figures above.
\clearpage
